%! Tex main = MarquezSalazarBrandon-PropuestaDeTesis.tex

\subsection{Objetivos}

\subsubsection{Objetivo general}

Implementar diversos NIM para la optimización de los
coeficientes de una función potencial en el marco de la
K-teoría dentro de la teoría de cuerdas, con el fin de
identificar y caracterizar una mayor cantidad de vacíos
estables.


\subsubsection{Objetivos específicos}

\begin{listaflor}
 \item Analizar los fundamentos teóricos de la K-teoría en teoría de
   cuerdas, con énfasis en la formulación de la función potencial
   y la noción de vacíos estables.
 \item Revisar y seleccionar métodos metaheurísticos adecuados
   (algoritmos genéticos, optimización por enjambre de
   partículas, algoritmos de estimación de distribución,
   colonia de hormigas, entre otros) para el problema
   planteado.
 \item Modelar el problema de optimización de los coeficientes de la
   función potencial como un problema computacional tratable
   mediante técnicas metaheurísticas.
 \item Implementar los métodos seleccionados y adaptar sus parámetros
   para su aplicación en el contexto del problema.
 \item Evaluar el desempeño de los métodos mediante métricas de
   convergencia, calidad de solución y capacidad para identificar
   vacíos estables.
 \item Comparar los resultados obtenidos con enfoques existentes o
   configuraciones base, a fin de determinar mejoras en la
   cantidad y calidad de los vacíos encontrados.
\end{listaflor}

